\documentclass[11pt,a4paper]{article}
\usepackage[T1]{fontenc}
\usepackage{lmodern}
\usepackage[margin=24mm]{geometry}
\usepackage{amsmath,amssymb,booktabs,longtable,array,hyperref}
\hypersetup{hidelinks}
\newcommand{\C}{\textbf{(c)}\ }
\newcommand{\D}{\textbf{(d)}\ }
\newcommand{\A}{\textbf{(a)}\ }
\newcommand{\file}[1]{\nolinkurl{#1}}
\newcommand{\alg}{\file{defect_bootstrap_algebra.py}}
\newcommand{\tra}{\file{defect_bootstrap_transmission.py}}
\newcommand{\sew}{\file{defect_bootstrap_sewing.py}}
\newcolumntype{L}[1]{>{\raggedright\arraybackslash}p{#1}}
\emergencystretch=3em
\title{Does a stable-vacuum quantum defect bootstrap\\remove the good-equation rank-two phase freedom?}
\author{Quantum Boussinesq solitons project --- PI: Gustav Delius}
\date{4 November 2026}
\begin{document}\maketitle
\begin{abstract}
\C \textbf{No removal follows from the specified constraints.} Their additional
rank on the exhibited antisymmetric coefficient is exactly zero. The cited
kernel is $F_{12}=c(P_1M_2-P_2M_1)$: rank two as a kernel, but only one
independent antisymmetric coefficient. Both ordered tensor coefficients
survive defect consistency before antisymmetry; antisymmetry itself leaves
one. Every nondegenerate real KP triangle has two light chords and one heavy
chord. Its proposed closed $2\leftrightarrow1$ cycle therefore has an unstable
bulk external leg, as well as unstable intermediate defect vacua. The
strictly stable-vacuum, light-particle absorption graph has no nontrivial edge.
Even allowing the entire triangle formally does not help: $F$ has zero
cycle sum and an explicit single-valued defect-transmission lift.

\C Two CZ defects compose on exact soliton backgrounds. Generically they
remain a two-parameter composite. A special reduction to a single oppositely
oriented CZ transmission factor exists on a KP triangle, but cannot have two
strictly stable input vacua. The same lift preserves that special reduction.
\D Existence of the quantum amplitudes, admissible resonance continuation,
residues, and the physical coefficient $c$ remain open. This is a derivation
of insufficiency of the proposed algebra, not a family of constructed quantum
field theories or a computation of the absolute physical phase.
\end{abstract}

\section{Scope, conventions, and audit precedence}
Labels are (a) literature proof, (b) literature heuristic, (c) project
calculation, and (d) unestablished physical identification or assumption.
All new checks are in the three scripts named above, under
\file{phase2/scripts/}; their JSON files record precision, errors and runtime.
No legacy script is run or imported. No git operation or renormalisation
choice is made.

\C The controlling sources are \file{audit_good_twobody.tex},
\file{audit_defq.tex}, \file{audit_defb_boundary.tex},
\file{audit_resonant_channels.tex}, and \file{errata.tex}. Their package
partners, the master report and README were read. For the classical
construction, \file{corrigan-zambon/verify_defect_solutions.py} and the
three-soliton sewing implementation were inspected; the new checker derives
the variational sewing equations afresh. The confirmed Casimir charges,
envelope theorem and canonical-charge holonomy do not supply normalized
quantum absorption vertices. We do not use the refuted counterterm
uniqueness, zero quantum parameter shift, or selection of $\lambda=0$.

\C The later \file{twobody_toda.tex} and \file{paper/review_draft3.tex}
correct the four-term Krein source printed in the frame audit. In that
subtracted frame it is not simply the previously advertised rank-one
product. Uniqueness of the specified Toda solution uses its corrected
source \emph{and axis data}. It is neither a consequence of diagonal
factorisation nor a good-equation input. The review also restricts the
minor test to a specified finite-rank ansatz. We do not infer that $F$ is
an exhaustive classification of the continuum ambiguity.

The fixed definitions are
\[
 u_{tt}=u_{xx}-u_{xxxx}-3(u^2)_{xx},\quad u=\phi_x=2(\log\tau)_{xx},
 \qquad \Omega=\int d\pi\wedge d\phi,
\]
\[
 \mathcal L=\tfrac12\phi_t^2-\tfrac12\phi_x^2+\phi_x^3
               -\tfrac12\phi_{xx}^2,\qquad \pi=\phi_t.
\]
\A The CZ source supplies the classical defect action and one-/two-soliton
transmission solutions \cite[Secs.~2, 3.4, 5.1--5.2]{CZ}.
\C The dictionary is $U_{\rm CZ}=-3\phi$, $\mathcal L_{\rm CZ}=9\mathcal L$.
We use $s=\sigma_{\rm vac}$ for scattering and $\eta s$ for the spectrum;
$Q=-s/(2\sqrt3)$, $\eta=\pm1$. This corrects the old one-loop note's
identification $s=\eta\sigma_{\rm vac}$. Below $q$ is phonon wavenumber,
and $Q$ is a defect root.

\section{The actual freedom and its frame}
\paragraph{Definitions and checks \C\ (\alg, \tra).}
For a soliton label $j=(\alpha_j,\epsilon_j)$,
\begin{gather}
 k_j=\sin\alpha_j,\quad v_j=\epsilon_j\cos\alpha_j,\quad
 M_j=2k_j,\quad \Pi_j=-2k_jv_j,\quad P_j=\tfrac23 k_j^3v_j,\quad
 \ell_j=-\epsilon_j\cos3\alpha_j,\label{data}\\
 a_j=\tfrac12(v_j/\sqrt3+k_j),\quad
 b_j=\tfrac12(v_j/\sqrt3-k_j),\qquad
 a_j^2+a_jb_j+b_j^2=\tfrac14,\label{conic}\\
 A_{12}=\frac{(a_1-a_2)(b_1-b_2)}{(a_1-b_2)(b_1-a_2)},\qquad
 \delta_0(1,2)=\tfrac29(\ell_1-\ell_2)\log A_{12}
                  +\tfrac23M_1M_2(v_1-v_2).\label{bulkclass}
\end{gather}
For regular two-body scattering we restrict to $0<\alpha_j<\pi/3$,
$A_{12}>0$, and unequal velocities. Labels stay on worldlines; reversing the
complete ordered pair reverses the generator. Coordinates are the
chamber-reversing Darboux coordinates of \file{good_twobody.tex}.
Relative to the fixed incoming chamber the audited subtraction is
\begin{equation}
 S_{\rm cb}=\tfrac12(\Pi_1M_2+\Pi_2M_1),\qquad
 \delta_0^{\rm raw}=\delta_0+S_{\rm cb}.\label{frame}
\end{equation}
This specifies the frame; a raw fixed-chamber phase is not thereby
antisymmetric. Below all amplitudes use this same convention and explicit seams.

\C Write the conditional expansion
\begin{equation}
 S(1,2)=\exp\{i\delta_0/\hbar+i[B_{12}^{\rm known}(\lambda)
                      +U_{12}+F_{12}]+O(\hbar)\}.\label{S}
\end{equation}
For definiteness the explicitly computable representative is
\begin{equation}
 B_{12}^{\rm known}(\lambda)=J_{12}
             +\sum_j\Delta X'_j p_{1,m}(j)+\lambda K_\lambda,\qquad
 J_{12}=\frac1{8\pi}\int(\vartheta_1\vartheta'_2-
                             \vartheta_2\vartheta'_1)\,dq.\label{knownbulk}
\end{equation}
The integral uses the half-line phases defined below, and $p_{1,m}$ is
\eqref{pm}. This is a conditional trace representative, not the completed
physical determinant. $U_{12}$ includes the uncomputed projected determinant,
state map, normalization and any other finite couplings. The new checker
independently integrates both ordered terms in overtaking and head-on
channels and verifies their endpoint sum. When isolating $F$, these other
unknown terms are held fixed; no claim of uniqueness of this decomposition
is made.
\D Identification with a physical elastic amplitude is an assumption.
\C The displayed counterexample in the cited package is precisely
\begin{equation}
 F_{12}=c(\lambda)(P_1M_2-P_2M_1).\label{F}
\end{equation}
Its two \emph{one-body} basis functions are $P$ and $M$. More generally,
$F_{12}=A P_1M_2+B M_1P_2$ obeys exchange antisymmetry iff $A+B=0$.
Thus $\wedge^2\operatorname{span}\{P,M\}$ is one-dimensional; there are
not two independent antisymmetric directions in the cited $F$.
The mixed $3\times3$ minors vanish identically, whereas a mixed $2\times2$
minor can be nonzero. Using mixed minors avoids the automatic vanishing
of an odd-dimensional skew determinant. This does not assert that every
continuum phase must have rank two.

\paragraph{Counterterm bookkeeping \C\ (\tra, \alg).}
For $\Delta H=\hbar\lambda\int u^2dx$, the computed contribution
to $B^{\rm known}$ in this frame is $\lambda K_\lambda$, where
\begin{align}
 \gamma_j&=-\frac{3\cos2\alpha_j}{\sin\alpha_j\sin3\alpha_j},\nonumber\\
 \Delta X'_1&=-\frac{\log A_{12}}{k_1}+\gamma_1M_2,
 &\Delta X'_2&=\frac{\log A_{12}}{k_2}-\gamma_2M_1,\nonumber\\
 K_\lambda&=3\delta_0+4\sum_jP_j\Delta X'_j,
 &\partial_{P_j}\delta_0&=-\Delta X'_j.\label{K}
\end{align}
At $(\alpha_1,\alpha_2)=(.3,.6)$,
\[K_\lambda^{++}=0.05008140486841890958\ldots,\qquad
K_\lambda^{+-}=-0.22517393868334175848\ldots.\]
The coefficient $c$ is independent of $\lambda$ as an \emph{unknown in the
constraint problem}; it is not proved physically scheme independent and
may be $c(\lambda)$. $K_\lambda$ is not $F$, and the parameter in the older
counterexample $A_\lambda$ is not this bulk counterterm.
An $O(\hbar)$ shift of the charge arguments of $F$ changes the exponent only
at $O(\hbar)$, so its $O(1)$ constraint uses classical charges at any fixed
finite $\lambda$.

\section{Stable-vacuum transmission: independently available input}
\paragraph{Classical part \C\ (\tra, \sew).}
Let $\theta=\operatorname{artanh}v$ if a signed rapidity is desired;
$k=\operatorname{sech}\theta$. We also denote the amplitude by
$T_Q(\alpha,\epsilon)$. No relativistic rapidity-difference assumption is made.
The left-to-right tau multiplier, directional delay, and classical generator are
\begin{align}
 z_{Q,\eta}&=\left(\frac{b-Q}{a-Q}\right)^\eta,&
 z_{t,Q}&=z_{Q,\eta}^{\epsilon},&
 \Delta x_Q&=-\frac{\log|z_{t,Q}|}{k},\label{z}\\
 \delta_Q^0&=\tfrac29[\ell-\ell(Q)]\log|z_{t,Q}|
                  -\tfrac23(\Pi M_Q^{\rm eff}-\Pi_Q^{\rm eff}M),\label{deltaD}\\
 \ell(Q)&=3\sqrt3(4Q^3-Q),&
 M_Q^{\rm eff}&=\epsilon\eta Q,&
 \Pi_Q^{\rm eff}&=\epsilon\eta(6Q^2-1)/\sqrt3,\nonumber\\
 \partial_P\delta_Q^0&=-\Delta x_Q.\label{deltagrad}
\end{align}
The effective coefficients are not intrinsic vacuum charges. Physical elastic
transmission is restricted to nonsingular $z_{Q,\eta}>0$; stable vacuum does
not ensure regular transmission at every light label. The real logarithmic
continuation elsewhere is algebraic only. The tree integration constant is
held fixed as in this audited generator.

\paragraph{Continuum input and its branch \C\ (\tra).}
Define phonon roots and the root factor
\begin{equation}
 a_q=\frac{\operatorname{sgn}(q)\sqrt{1+q^2}}{2\sqrt3}+\frac{iq}{2},
 \quad b_q=\overline{a_q},\qquad
 Z(q;Q)=\frac{b_q-Q}{a_q-Q},\quad z_{\rm ph}=Z(q;Q)^\eta.
 \label{Z}
\end{equation}
The latter has limit $e^{-2\pi i\eta/3}$ at \emph{both} infinities.
The soliton--phonon factor is $Z(q;b)/Z(q;a)$.
Use the soliton phase continuous on each open half-line and zero at infinity:
\begin{align}
 \vartheta_S(q)=-\sum_{r=1}^3w_r\Big[\operatorname{sgn}(q)
 \arctan\frac{\sin\varphi_r}{|q|}&\nonumber\\[-1ex]
 {}+\epsilon\arctan\frac{\sin\varphi_r\cos\varphi_r}
 {|q|\sqrt{1+q^2}+q^2+\sin^2\varphi_r}\Big],&\label{thetaS}\\
 (\varphi_1,\varphi_2,\varphi_3)=(\alpha+2\pi/3,\alpha-2\pi/3,\alpha),
 \quad(w_1,w_2,w_3)=(1,1,-2).&\nonumber
\end{align}
Let $\vartheta_Q^0$ be the continuous argument of $Z^\eta$ starting from
zero separately at $q=0^+$ and $q=0^-$. With ordinary half-line integrals,
\begin{equation}
 I_{SQ}^0=\frac1{4\pi}\int_{\mathbb R\setminus\{0\}}
              \vartheta_S\,\partial_q\vartheta_Q^0\,dq
       =-\frac1{4\pi}\int_{\mathbb R\setminus\{0\}}
              \vartheta_Q^0\,\partial_q\vartheta_S\,dq.
 \label{I}
\end{equation}
No delta function is inserted at zero. Integration-by-parts endpoints vanish
in this branch convention. This is the direct continuum integral, not an
identification with a different phase reference.

\C At $(\alpha,\epsilon,Q,\eta)=(.35,+1,.5,+1)$,
\begin{equation}
 I_{SQ}^0=0.7622851906696779106824093224\ldots .\label{bench}
\end{equation}
The boundary audit's infinity-reference integration-by-parts table instead has
$I_{SQ}^{\infty}=I_{SQ}^0-2\pi/3
=-1.3321099117235175816260195998\ldots$ at this point.
Here $\vartheta_S(0^+)=2\pi$, $\vartheta_S(0^-)=0$; the defect reference
shifts are $-4\pi/3$ and $+2\pi/3$ on the positive and negative half-lines.
This checked endpoint correction is not an absolute physical discrepancy.
Changing representative requires compensating the unknown amplitude term.

\paragraph{The honest $O(\hbar^0)$ amplitude \C/\D\ (\tra).}
For a strictly stable vacuum, $\eta Q>1/(2\sqrt3)$, write
\begin{equation}
 \boxed{T_Q(\theta)=\exp\left\{\frac{i}{\hbar}\delta_Q^0
       +i\left[I_{SQ}^0+\Delta x_Q\,p_1(\theta;\lambda)
                    +W_Q(\theta;\lambda)\right]+O(\hbar)\right\}.}
 \label{T}
\end{equation}
The Taylor term has a \emph{plus} sign since $P^{\rm phys}=P+\hbar p_1$
and \eqref{deltagrad}. This is the boundary audit's corrected sign;
a new fixed-$P$ finite difference verifies it. A momentum representative
must have the matching label-frame determinant. For example the mode trace is
\begin{equation}
 p_{1,m}=-\frac{\epsilon}{2\pi}\sum_rw_r\sin\varphi_r
                                  \log|\tan(\varphi_r/2)|.\label{pm}
\end{equation}
The canonical momentum counterterm contribution is $2\lambda vk$; it is zero
in the fixed-profile mode momentum. A change $p_1\mapsto p_1+\Delta p_1$
requires $W_Q\mapsto W_Q-\Delta x_Q\Delta p_1$ for the same amplitude.

\D $W_Q$ is not set to zero. It contains the missing boundary-compatible
projected determinant, finite boundary terms, root/parameter shifts acting on
$\delta_Q^0$, the label-frame bulk deformation and phase normalization.
The Casimir/envelope formulas do not determine it. For this trace convention
only, $I^0+\Delta x_Qp_{1,m}$ at \eqref{bench} is
$2.2646110601221966897621495006\ldots$ \C; it has no error bound as an
approximation to a full physical phase \D. Neither $I^0$ nor $\delta_Q^0$
is defined using the unknown bulk phase.

\section{The proposed consistency equations}
\subsection{Two solitons and a stable defect}
\D Suppose the light elastic sector is scalar and diagonal. Products act on
ordered incoming worldlines with the rightmost encounter first. Require
\begin{equation}
 S(1,2)T_Q(1)T_Q(2)=T_Q(2)T_Q(1)S(1,2).\label{RTT}
\end{equation}
\C In the same Darboux/seam convention on both sides, factors commute.
Its logarithm vanishes for arbitrary $F$ and $W_Q$, so its constraint rank
is zero. This is scalar algebra, not a field-theory factorisation theorem.
Changing chamber convention would require explicit intertwiners on both
sides. The nonlinear and phonon checks below test the classical field input
independently rather than counting commutativity as a field calculation.

\subsection{Absorption poles, intertwiners, and the closed cycle}
\D A quantum absorption bootstrap needs a meromorphic continuation, an
admissible state $D_j$, and a nonzero normalized vertex
$\Gamma_i:D_i\otimes S_i\to D_j$. In an analytic coordinate $\zeta_i$, a
conventional pole requirement would have the form
\begin{equation}
 T_i(\zeta_i)\sim\frac{i\,\Gamma_i^\dagger\Gamma_i}
                             {\zeta_i-\zeta_i^*}+\text{regular terms}.
 \label{pole}
\end{equation}
The residue sign, coordinate Jacobian and sheet need spectral normalization;
\eqref{pole} is an assumption, not a measured residue. No audited input
supplies them. A classical zero/pole of $z_Q$ is a divergent delay, not a
simple quantum pole of $T_Q$: the coefficient of its logarithm in
\eqref{deltaD} vanishes at absorption. The loop expansion need not be uniform
there. We do not manufacture a residue from that expansion.

\C The independently checked tree identity fixes the required seam.
For $Q_i\to Q_j$, $S_i=\{Q_i,Q_j\}$ with positive $k_i$,
\begin{equation}
 \rho_i=\eta\operatorname{sgn}(Q_i-Q_j),\quad
 \sigma_{ti}=\epsilon_t\rho_i,\qquad
 \delta^0_{Q_j}(t)-\delta^0_{Q_i}(t)
     =\sigma_{ti}[\delta_0(t,i)+M_i\Pi_t].\label{seam}
\end{equation}
Both test directions and defect orientations are checked in \tra.
\D If a vertex exists the scalar intertwining condition is
\begin{equation}
 T_{Q_j}(t)\Gamma_i=C_{ti}S(t,i)^{\sigma_{ti}}\Gamma_i T_{Q_i}(t),\qquad
 C_{ti}=\exp\{i\sigma_{ti}M_i\Pi_t/\hbar+i\chi_{ti}+O(\hbar)\}.
 \label{intertwine}
\end{equation}
$S^{-1}$ means the inverse oriented ordered amplitude. The uncomputed
$\chi_{ti}$ includes the quantum seam and shifted-kinematics terms.
Removing it without computing it would not be a fixed-frame bootstrap.
For emission use the inverse intertwiner, not a second stable absorption state.

\C Conditional on nonzero vertices, the $O(1)$ logarithm gives
\begin{equation}
 t_{Q_j}(t)-t_{Q_i}(t)=\sigma_{ti}b(t,i)+\chi_{ti},\qquad
 b=B^{\rm known}+U+F,\quad t_Q=I^0+\Delta x_Qp_1+W_Q.\label{edgeeq}
\end{equation}
The cycle telescopes to the necessary condition
\begin{equation}
 \sum_{i=1}^3[\sigma_{ti}b(t,i)+\chi_{ti}]=0\pmod{2\pi}.\label{cycle}
\end{equation}
Equivalently a spectator commutes through a nonzero closed fusion vertex.
For homogeneous infinitesimal changes the branch integer is fixed, so the
sum is exactly zero. This is a conditional algebraic test even if the
physical vertices do not exist.

\paragraph{Which states occur \C\ (\alg; controlling instability audit).}
For $0<\vartheta<\pi/3$, parametrize the ordered KP roots by
\begin{equation}
 Q_1=\frac{\cos\vartheta}{2\sqrt3}+\frac{\sin\vartheta}{2},\quad
 Q_2=\frac{\cos\vartheta}{2\sqrt3}-\frac{\sin\vartheta}{2},\quad
 Q_3=-\frac{\cos\vartheta}{\sqrt3}.\label{roots}
\end{equation}
They obey $Q_1>Q_2>Q_3$ and share $Q^3-Q/4$.
The walk $1\to2\to3\to1$ has positive-mass chords
\begin{align}
 k_1&=\sin\vartheta,&v_1&=\cos\vartheta,\nonumber\\
 k_2&=\sin(\pi/3-\vartheta),&v_2&=-\cos(\pi/3-\vartheta),\nonumber\\
 k_3&=\sin(\pi/3+\vartheta)=k_1+k_2,&
 v_3&=-\cos(\pi/3+\vartheta).\label{chords}
\end{align}
Thus $\alpha_1=\vartheta$ and $\alpha_2=\pi/3-\vartheta$ are light, whereas
\begin{equation}
 \alpha_3=\min(\pi/3+\vartheta,2\pi/3-\vartheta)>\pi/3.\label{heavy}
\end{equation}
At $\vartheta=\pi/6$, $\alpha_3=\pi/2$ is static, still heavy and unstable.
The instability audit constructs a localized growing mode of this chord;
its positive rate is $\sqrt3k_1k_2=\sqrt3(1/4-v_3^2)$, with decaying tails
since $k_3>\sqrt3|v_3|$. These kinematic identities are rechecked here; the
spectral proof is inherited from that audit. The bar on $\bar S_3$ denotes
the emitted outer chord, not a light antiparticle with negative mass.
Relabelling can move the heavy chord to $S_1$ or $S_2$, but cannot remove it.

\C The stable root is $Q_1$ for $\eta=+1$ and $Q_3$ for $\eta=-1$:
$\eta Q>1/(2\sqrt3)$ is the strictly stable half-line and $12Q^2-1>0$
makes the common cubic injective there. The middle root lies between the
thresholds. Direct admissible-root filtering of the audited sewing quadratics
confirms the count. Marginal endpoints are excluded. For negative orientation
reverse/relabel the walk to start at the stable vertex. Neither orientation
has stable intermediate vacua.

\C Deleting the heavy edge leaves the path $Q_1-Q_2-Q_3$: every closed walk
reduces to backtracks. Deleting unstable vertices as asymptotic states leaves
one isolated vertex. Absorption preserves the KP cubic and cannot move to
another triangle within this graph. \D A light-only diagonal transmission
sector can still be postulated and tested by \eqref{RTT}; a nontrivial
light-only bound-state bootstrap is not obtained from these edges. Virtual
unstable vacua and heavy resonances need a separate analytic theory. This
does not exclude quantum states or composite vertices absent from the inputs.

\subsection{Two defects in series and exceptional reduction}
\paragraph{Classical composition \A/\C\ (\sew, \alg, \tra).}
The CZ two-soliton solution has $\tau(E_1,E_2)=1+E_1+E_2+A_{12}E_1E_2$
and implements a defect by $E_j\mapsto z_jE_j$ \cite[Sec.~5.2]{CZ}.
Two defects at $x=0,L$ therefore give
\begin{equation}
 \tau_L=\tau(E_j),\quad \tau_M=\tau(z_{1j}E_j),\quad
 \tau_R=\tau(z_{2j}z_{1j}E_j).\label{series}
\end{equation}
The middle field is retained. \sew\ varies the local CZ action and checks
all four sewing equations at both locations, four orientation pairs and three
times. It also checks the bulk Hirota identity and the linearized
two-soliton-plus-phonon mode at both defects. Sampled defects are strictly
stable and transmissions regular. Their scalar tau maps commute. The
asymptotic phonon transmission is
\begin{equation}
 T_{\rm ph}^{\rm series}(q)=\prod_j\frac{Z(q;b_j)}{Z(q;a_j)}
                   Z(q;Q_1)^{\eta_1}Z(q;Q_2)^{\eta_2}.\label{phseries}
\end{equation}
This explicit mode has no reflected asymptotic component. Completeness and
stability of all localized modes of the two-defect system do not follow.

\C Generically this is not a single CZ defect with an added parameter.
The high-frequency phase requires $\eta_1=\eta_2=\eta$ and
$\eta_{\rm eff}=-\eta$ for an elementary reduction. Opposite orientations
give limit $1$, unlike an elementary factor. For equal orientations,
cross-multiplication on the conic \eqref{conic} gives
\begin{equation}
 z_{Q_1,\eta}z_{Q_2,\eta}=z_{Q_3,-\eta}
 \quad\Longleftrightarrow\quad Q_3=-Q_1-Q_2,\quad
 Q_1^2+Q_1Q_2+Q_2^2=\tfrac14.\label{special}
\end{equation}
Indeed the difference of the two cubic products divided by $b-a$ reduces to
\begin{equation}
 \tfrac14-(Q_1+Q_2+Q_3)(a+b)+Q_1Q_2+Q_1Q_3+Q_2Q_3.\label{poly}
\end{equation}
It vanishes for all test solitons exactly under \eqref{special}; the same
identity holds for phonons. This reduces the classical scattering/tau map,
not necessarily the local action after eliminating the middle field.

\C Two strictly stable equal-orientation inputs have
$\eta Q_1=1/(2\sqrt3)+u$ and $\eta Q_2=1/(2\sqrt3)+w$, $u,w>0$. Then
\begin{equation}
 Q_1^2+Q_1Q_2+Q_2^2-\tfrac14
 =u^2+uw+w^2+\tfrac{\sqrt3}{2}(u+w)>0.\label{nostablefusion}
\end{equation}
Thus the special reduction is unavailable. The inverse pair at equal root
and opposite orientations gives the identity map, but both vacua cannot be
strictly stable. Threshold/repeated-root limits do not yield a generic
strictly stable fusion law.

\paragraph{At one loop \C/\D\ (\tra, \sew).}
For separated encounters the proposed scalar factorisation is
\begin{equation}
 T_{D_2D_1}(t)=T_{D_2}(t)T_{D_1}(t),\qquad
 \delta_{D_2D_1}^0=\delta_{D_1}^0+\delta_{D_2}^0,\quad
 t_{D_2D_1}=t_{D_1}+t_{D_2}.\label{quantseries}
\end{equation}
Free propagation between reference planes is factored out. Classical and
continuum-trace additivity follow from \eqref{series}, logarithmic derivatives
of \eqref{phseries}, and linearity of \eqref{I}. Physical quantum factorisation
is still an assumption. A contact limit can have new terms or retained internal
degrees of freedom, and cannot be declared an elementary CZ parameter without
an action calculation.

\C Even on the formal special locus the classical generators add with the
chosen constants. Phase derivatives and the zero-referenced integral $I^0$
also add, since the three-root product is one and branches start at $0^\pm$.
Frame constants must be transformed consistently. The next section's homogeneous
freedom preserves this exceptional reduction too.

\section{Insertion of the freedom and the constraint rank}
\paragraph{Independent charge primitives \C\ (\alg).}
Define polynomial vertex functions
\begin{equation}
 m(Q)=2Q,\qquad p(Q)=2\sqrt3Q^4-\frac{Q^2}{\sqrt3}-\frac{\sqrt3}{36}.
 \label{potentials}
\end{equation}
Polynomial reduction modulo $a^2+ab+b^2=1/4$ gives
\begin{equation}
 M=m(a)-m(b),\qquad P=p(a)-p(b).\label{edgecharges}
\end{equation}
This derives primitives from the one-body charges, without a two-body phase
or quantum defect assumption. On each KP triangle,
\begin{equation}
 \sum_i\rho_iM_i=\sum_i\rho_iP_i=0,\qquad
 \sum_{i=1}^3m(Q_i)=\sum_{i=1}^3p(Q_i)=0.\label{sums}
\end{equation}
The constant in $p$ does not affect edge differences. It centers the
three-vertex sum using $\sum Q_i^2=1/2$ and $\sum Q_i^4=1/8$; this is a
witness for composition, not a chosen physical counterterm.

\C Inserting $F_{ti}=A P_tM_i+B M_tP_i$ into the cycle gives
\begin{equation}
 \sum_i\sigma_{ti}F_{ti}=\epsilon_t\left[
 A P_t\sum_i\rho_iM_i+B M_t\sum_i\rho_iP_i\right]=0.\label{Fcycle}
\end{equation}
Each coefficient cancels separately, without setting $c$ or $\lambda$ to zero.

\paragraph{A global null direction \C\ (\alg, \tra).}
When comparing two formal solutions, hold the known transmission, quantum
seam and other nuisance terms fixed. The single-valued change
\begin{equation}
 \Delta W_Q(t)=-\epsilon_t\eta[A P_t m(Q)+B M_t p(Q)]\label{lift}
\end{equation}
satisfies every homogeneous edge equation:
\begin{equation}
 \Delta W_{Q_j}(t)-\Delta W_{Q_i}(t)=\sigma_{ti}F_{ti}.\label{liftedge}
\end{equation}
For $A=c,B=-c$ it becomes
\begin{equation}
 \boxed{\Delta W_Q(t)=\epsilon_t\eta\,c[M_t p(Q)-P_t m(Q)].}\label{liftc}
\end{equation}
It is real and smooth on regular light domains, vanishes as the test soliton
disappears and changes no leading classical or phonon factor. Its local
analytic continuation is regular away from existing charge-coordinate branch
endpoints, so it does not introduce poles itself; unspecified residues can
be rescaled. It is not a rephasing of the defect vacuum alone, because it
depends on the incident soliton. \D No claim is made that a local integrable
quantum boundary operator generates it. The statement is that the given
algebra and available input do not exclude it.

\C For two defects add their \eqref{lift}; on the special locus,
\begin{equation}
 \Delta W_{Q_1,\eta}(t)+\Delta W_{Q_2,\eta}(t)=\Delta W_{Q_3,-\eta}(t).
 \label{liftfusion}
\end{equation}
This preserves exceptional reduction as well as the stable separated product.
The lift is derived as a null vector \emph{after} specifying the independent
known part of $T$. It is not used to define that known part, and cannot
be advertised as a derivation of the physical $S$ from $T$.

\paragraph{Explicit rank computation \C\ (\alg).}
For a fixed spectator collect $H_i=\Delta W_{Q_i}(t)$ into a column. The
three edge equations are
\begin{equation}
 DH=E\begin{pmatrix}A\\B\end{pmatrix},\quad
 D=\begin{pmatrix}-1&1&0\\0&-1&1\\1&0&-1\end{pmatrix},\quad
 E_i=\sigma_{ti}(P_tM_i,M_tP_i).\label{rankeq}
\end{equation}
The left null vector $(1,1,1)$ annihilates $E$, so
\begin{equation}
 \boxed{\operatorname{rank}[D\mid E]-\operatorname{rank}D=2-2=0.}\label{rank0}
\end{equation}
The script checks the exact matrix at the nondegenerate symmetric triangle
and proves the global lift for arbitrary triangle parameter; no numerical
singular-value threshold is used. Writing the lift more generally as
$\epsilon_t\eta[dP_tm(Q)+eM_tp(Q)]$, edge equations reduce to $d=-A$,
$e=-B$: rank two in four variables, zero remaining rank on $(A,B)$ after
eliminating $(d,e)$. The separate row $A+B=0$ leaves the single coefficient $c$.

\begin{center}\small
\begin{tabular}{@{}L{8.5cm}cc@{}}\toprule
Constraint & New rank on $(A,B)$ & New rank on $c$\\\midrule
Scalar soliton/defect factorisation & 0 & 0\\
Strictly stable, light-only absorption graph & 0 & 0\\
Formal full triangle, transmission freedoms retained & 0 & 0\\
Closed cycle alone, even with fixed vertex phases & 0 & 0\\
Separated two-defect consistency & 0 & 0\\
Exceptional CZ reduction, with \eqref{lift} & 0 & 0\\
Exchange antisymmetry (already imposed) & 1 & ---\\\bottomrule
\end{tabular}\end{center}

\D These are homogeneous ranks of necessary conditional equations. We have
not proved that their inhomogeneous known terms admit a physical solution.
The confirmed canonical-charge holonomy and absorbed-leg Taylor remainder
obstruct naive unshifted-edge identifications. $F$ cannot repair a nonzero
cycle remainder because \eqref{Fcycle} is zero. Shifted quantum kinematics,
seams and vertex normalization need quantum sewing. Holonomy is neither
a condition $c=0$ nor a universal fusion no-go. Independently fixing $W_Q$
and physical residues, together with valid vertices, could make individual
edge equations informative; this is additional input, absent here. The null
direction is not evidence for an adjustable physical coupling.

\section{Verification, errors and failures recorded}
\begin{longtable}{@{}L{5cm}L{3.2cm}L{7cm}@{}}\toprule
Claim and label & Checker & Coverage\\\midrule\endhead
\C Kernel rank, coefficient count, charge primitives, cycle rank zero
& \alg & Exact polynomial reduction, mixed minors, incidence rank,
global lift and exceptional composition.\\
\C Light/light/heavy chords and stable roots & \alg & Rational triangle
identities, analytic inequalities, 80-digit geometry and admissible-root
filtering. Instability spectral proof inherited from its audit.\\
\C Generator, delay sign, seam and counterterm response & \tra & Independent
derivatives, both orientations/directions; fixed-$P$ finite difference.\\
\C Continuum trace and reference constant & \tra & Two half-line quadrature
forms and 60/80-digit replication; explicit endpoint correction.\\
\C Two-defect fields and phonons & \sew & Action-derived sewings: 96 nonlinear
and 32 linearized equations, both locations; 70 digits; bulk Hirota check.\\
\D Quantum poles, physical amplitudes, local quantum realization of the lift
& None claimed & Not established by the successful algebraic checks.\\
\bottomrule\end{longtable}

\paragraph{Error budget \C.}
The rank \eqref{rank0} is exact and has no numerical uncertainty. The largest
sampled charge/chord closure residual is $1.06\times10^{-81}$ at 80 digits
(assertion $10^{-65}$). The two integral forms differ by
$1.06\times10^{-81}$; 60/80-digit replication changes \eqref{bench} by
$2.12\times10^{-62}$ (assertion $10^{-55}$). These are deterministic agreement
diagnostics, not rigorous quadrature intervals. The boundary-audit table has
only twenty decimal places, so its comparison uses $6\times10^{-20}$ and
finds $2.61\times10^{-20}$. The fixed-$P$ central difference uses step
$10^{-20}$, error $2.37\times10^{-40}$ and assertion $10^{-35}$.
Nonlinear sewing residual is $6.68\times10^{-70}$, linearized residual
$3.16\times10^{-71}$, both below $10^{-52}$. No bound on an uncomputed
physical phase follows from these numbers.

\paragraph{Failures recorded.}
\begin{enumerate}
\item \C The interpretation of ``both directions'' of $F$ fails: its source
has one antisymmetric coefficient on two one-body basis functions. Both
ordered coefficients were retained before imposing antisymmetry.
\item \C The stable/light external three-edge absorption cycle fails its
admissible-state test: the outer chord is heavy. The exceptional two-CZ
reduction fails the two-strictly-stable-input test. These are negative results.
\item \C Inherited failed premises remain failures: selection of $\lambda=0$,
unique zero quantum parameter shift, the old transmission Taylor sign and
the uncorrected four-term Krein source. Their audit corrections are retained.
No residue is inferred from a classical delay pole.
\item \C Exploratory file discovery encountered protected temporary-directory
permission diagnostics and missing diagnostics for not-yet-created outputs.
No scientific check depended on these searches.
\item \C The first report-assembly command used a repository-relative path
while already in \file{phase2/}. The source was not created and the PDF build
failed with a missing-file diagnostic. The corrected path repaired this;
the failed build's temporary unprefixed log was removed immediately.
The next build exposed a delimiter spanning an alignment row and a long
inline number; the delimiter and line layout were corrected. Rendered-page
inspection then caught a stray delimiter fragment in the phase formula and
spaces lost by the URL-style command font; both were repaired. The first
formatting repair repeated the relative-path mistake and made no edit;
using an absolute source path fixed it.
\item \C All three initial scientific runs passed, as did the strengthened checks
of the displayed charge identities and continuum representatives. No assertion, tolerance
or precision was relaxed. No failed numerical comparison was hidden.
\end{enumerate}

\paragraph{Conclusion.}
\C The stated defect constraints impose zero additional rank on $c$.
The formal cycle is blind to $F$ because its two basis charges are chord
differences. The single-valued lift preserves formal edges and exceptional
two-defect reduction. The strictly stable/light sector has fewer bootstrap
relations. \D Determining $c(\lambda)$ still needs an independently derived
boundary-compatible quantum determinant/state map or a quantum boundary
theory with admissible poles and normalized vertices. Nothing here transfers
a Toda prescription to the good equation or promotes an unstable vacuum
to an asymptotic state.

\paragraph{Reproducibility and reading scope.}
Run \texttt{python3 }\file{phase2/scripts/defect_bootstrap_algebra.py} and the analogous
commands for \file{defect_bootstrap_transmission.py} and
\file{defect_bootstrap_sewing.py}. Each writes only its own prefixed JSON
and ends with \texttt{ALL CHECKS PASSED}. Initial measured runtimes were
3.8, 7.4 and 8.3 seconds. The strengthened transmission run took 26.3 seconds;
all final scientific checks together take under 40 seconds on this host.
The final PDF was built by two \file{pdflatex} passes from \file{phase2/};
all ten rendered pages were inspected.
The sole external reference's entry, abstract and relevant full-text passages
in sections 2, 3.4 and 5 were read online; the action is also checked
against the repository derivations. The entire article was not read.
No new FK attribution or unverified identifier is used.

\begin{thebibliography}{9}
\bibitem{CZ} E.~Corrigan and C.~Zambon,
\emph{Adding integrable defects to the Boussinesq equation},
\href{https://arxiv.org/html/2306.03772v1}{arXiv:2306.03772v1}.
Reading scope stated above.
\end{thebibliography}
\end{document}
